\RequirePackage{filecontents}
\begin{filecontents*}{moderncvcolorcustom.sty}
\NeedsTeXFormat{LaTeX2e}
\ProvidesPackage{moderncvcolorcustom}[2013/02/09 v1.3.0 modern curriculum vitae and letter color scheme: custom]
\definecolor{color0}{rgb}{0,0,0}% black
% \definecolor{color1}{rgb}{0.031, 0.569, 0.698}
\definecolor{color2}{rgb}{0.45,0.45,0.45}% dark grey
\definecolor{FhGreen}{RGB}{23,156,125}
%
\definecolor{color1}{RGB}{23,156,125}
\endinput
\end{filecontents*}


%%%%%%%%%%%%%%%%%%%%%%%%%%%%%%%%%%%%%%%%
%
% run xelatex
% bibtex on file
% bibtex on new aux file
% xelatex
% xelatex
%
%%%%%%%%%%%%%%%%%%%%%%%%%%%%%%%%%%%%%%%%

\documentclass[11pt, a4paper, sans]{moderncv}        
% possible options include:
% font size ('10pt', '11pt' and '12pt'), 
% paper size ('a4paper', 'letterpaper', 'a5paper', 'legalpaper', 'executivepaper' and 'landscape') 
% font family ('sans' and 'roman')

% moderncv themes
\moderncvstyle{classic}                             
% style options are 'casual' (default), 'classic', 'banking', 'oldstyle' and 'fancy'

\moderncvcolor{custom}   
                            
% color options 'black', 'blue' (default), 'burgundy', 'green', 'grey', 'orange', 'purple', 'red' and 'custom'


% \usepackage[notbib, notlof, notlot, nottoc]{tocbibind}

%% ───────────────────────────────────────────────────────────────────────────────

\usepackage{tabularx}

\RequirePackage{FraunhoferColors}

%% use Fraunhofer colors beamer structure elements:
%\setbeamercolor{structure}{fg=FhGreen}
%\setbeamercolor{alerted text}{fg=FhDarkOrange}
%\setbeamercolor{example text}{fg=FhLightTeal}
%
%% colorize itemize elements:
%\setbeamercolor{item}{fg=FhGreen}
%\setbeamercolor{subitem}{fg=FhSilver}
%\setbeamercolor{subsubitem}{fg=FhSilver}

%% ───────────────────────────────────────────────────────────────────────────────
%% 1) needed packages
\usepackage{xparse}                % for \RenewDocumentCommand
\usepackage[x11names]{xcolor}      % for \textcolor and colour mixing
\usepackage{calc}                  % for \heightof
%
%% 2) redefine \section
\makeatletter
%
%%\RenewDocumentCommand{\section}{s O{} m}{%
%%  % reset internals (moderncv chain)
%%  \tl@resetchain%
%%  % vertical spacing before
%%  \par\addvspace{2.5ex}%
%%  % hyperref anchor
%%  \phantomsection
%%  % add to TOC
%%  \addcontentsline{toc}{section}{#3}%
%%  % the title itself: 
%%  %  · \sectionstyle is the font/size wrapper from moderncv
%%  %  · first \textcolor paints the rule in “awesome!55!white”
%%  %  · then the actual title text #3 follows
%%  \strut
%%  \sectionstyle{%
%%    \textcolor{awesome!55!white}{%
%%      #2\relax            % optional star‐prefix
%%      \rule{.75ex}{0pt}   % horizontal gap
%%      \rule{1pt}{\heightof{#3}}% vertical rule
%%    }%
%%    \rule{.75ex}{0pt}     % gap between rule and text
%%    #3                    % the actual section title
%%  }%
%%  % spacing after
%%  \par\nobreak\addvspace{1ex}\@afterheading
%%}
%
\definecolor{sectionrulecolor}{rgb}{0.031, 0.569, 0.698}

% 3) replace \section
\makeatletter
\renewcommand*{\section}[1]{%
  % reset moderncv internals
  
  %\tl@resetchain%

  % vertical space before
  \par\addvspace{2.5ex}%
  % hyperref anchor
  \phantomsection%
  % TOC entry
  \addcontentsline{toc}{section}{#1}%
  % the heading itself
  \strut
  \sectionstyle{%
    %  · draw the rule in our colour
    \textcolor{sectionrulecolor}{%
      %\rule{.75ex}{0pt}% horiz gap    
      %\rule[\dimexpr\heightof{#1}/4\relax]{20pt}{\dimexpr\heightof{#1}/2\relax}
      \rule[3pt]{19pt}{2pt}% vertical rule
    }%
    %  · gap then the title text
    \rule{.75ex}{0pt}#1%
  }%
  % spacing after
  \par\nobreak\addvspace{1ex}\@afterheading%
}
\makeatother
%% ───────────────────────────────────────────────────────────────────────────────

\definecolor{custom}{rgb}{0.031, 0.569, 0.698}
\definecolor{awesome}{rgb}{0.031, 0.569, 0.698}

\definecolor{custom}{RGB}{23,156,125}
\definecolor{awesome}{RGB}{23,156,125}

\definecolor{color0}{rgb}{0,0,0}% black
%\definecolor{color1}{rgb}{0.031, 0.569, 0.698}
\definecolor{color2}{rgb}{0.45,0.45,0.45}% dark grey

\definecolor{color2}{rgb}{0.45,0.45,0.45}% dark grey

\definecolor{sectionrulecolor}{RGB}{8,145,178}



\definecolor{blueXIIdark}{cmyk}{1,.8,.30,.05}
\definecolor{blueXIIlight}{cmyk}{.0,.30,1,.00}


%-------------------------------------------------
\usepackage{amsmath}
\usepackage{amsfonts}
\usepackage{amssymb}


\usepackage{tikz}
\usetikzlibrary{calc}
\usetikzlibrary{fit}
\usepackage{graphicx}

% GanttHeader setups some parameters for the rest of the diagram
% #1 Width of the diagram
% #2 Width of the space reserved for task numbers
% #3 Width of the space reserved for task names
% #4 Number of months in the diagram
% In addition to these parameters, the layout of the diagram is influenced
% by keys defined below, such as y, which changes the vertical scale
\def\GanttHeader#1#2#3#4{%
 \pgfmathparse{(#1-#2-#3)/#4}
 \tikzset{y=7mm, task number/.style={left, font=\bfseries},
     task description/.style={text width=#3,  right, draw=none,
           font=\sffamily, xshift=#2,
           minimum height=2em},
     gantt bar/.style={draw=black, fill=FhGreen!30},
     help lines/.style={draw=black!30, dashed},
     x=\pgfmathresult pt
     }
  \def\totalmonths{#4}
  \node (Header) [task description] at (0,0) {\textbf{\large Tasks}};
  \begin{scope}[shift=($(Header.south east)$)]
    \foreach \x in {1,...,#4}
      \node[above] at (\x,0) {\footnotesize\x};
 \end{scope}
}

% This macro adds a task to the diagram
% #1 Number of the task
% #2 Task's name
% #3 Starting date of the task (month's number, can be non-integer)
% #4 Task's duration in months (can be non-integer)
\def\Task#1#2#3#4{%
\node[task number] at ($(Header.west) + (0, -#1)$) {#1};
\node[task description] at (0,-#1) {#2};
\begin{scope}[shift=($(Header.south east)$)]
  \draw (0,-#1) rectangle +(\totalmonths, 1);
  \foreach \x in {1,...,\totalmonths}
    \draw[help lines] (\x,-#1) -- +(0,1);
  \filldraw[gantt bar] ($(#3, -#1+0.2)$) rectangle +(#4,0.6);
\end{scope}
}
%-------------------------------------------------












\renewcommand{\familydefault}{\sfdefault}         
% to set the default font; use '\sfdefault' for the default sans serif font, '\rmdefault' for the default roman one, or any tex font name

%\nopagenumbers{}                                  
% uncomment to suppress automatic page numbering for CVs longer than one page

\usepackage{fancyhdr}
% \usepackage{hyperref}

% adjust the page margins
\usepackage{geometry}

% footer details
\usepackage{lastpage}

% Configure page margins with geometry
\geometry{left=2.4cm, top=2.8cm, right=2.4cm, bottom=2.8cm, footskip=.5cm}

% Specify the location of the included fonts
%\fontdir[/usr/share/texlive/texmf-dist/fonts/]

% Color for highlights
% Awesome Colors: awesome-emerald, awesome-skyblue, awesome-red, awesome-pink, awesome-orange, awesome-nephritis, awesome-concrete, awesome-darknight


\setlength{\footskip}{36.00005pt}                 
% depending on the amount of information in the footer, you need to change this value. comment this line out and set it to the size given in the warning

%\setlength{\hintscolumnwidth}{3cm}                
% if you want to change the width of the column with the dates

%\setlength{\makecvheadnamewidth}{10cm}            
% for the 'classic' style, if you want to force the width allocated to your name and avoid line breaks. be careful though, the length is normally calculated to avoid any overlap with your personal info; use this at your own typographical risks...

% font loading
% for luatex and xetex, do not use inputenc and fontenc
% see https://tex.stackexchange.com/a/496643
\ifxetexorluatex
  \usepackage{fontspec}
  \usepackage{unicode-math}
  \defaultfontfeatures{Ligatures=TeX}
  \setmainfont{Latin Modern Roman}
  \setsansfont{Latin Modern Sans}
  \setmonofont{Latin Modern Mono}
  \setmathfont{Latin Modern Math} 
\else
  \usepackage[T1]{fontenc}
  \usepackage{lmodern}
\fi

\usepackage{academicons}
\RequirePackage[default,opentype]{sourcesanspro}

%-------------------------------------------------------------------------------
%                Configuration for fonts
%-------------------------------------------------------------------------------
% Set the FontAwesome font to be up-to-date.
%\newfontfamily\FA[Path=\@fontdir]{FontAwesome}
% Set font for header (default is Roboto)
\newfontfamily\headerfont[
  %Path=\@fontdir,
  UprightFont=*-Regular,
  ItalicFont=*-Italic,
  BoldFont=*-Bold,
  BoldItalicFont=*-BoldItalic,
]{Roboto}

\newfontfamily\headerfontlight[
  %Path=\@fontdir,
  UprightFont=*-Thin,
  ItalicFont=*-ThinItalic,
  BoldFont=*-Medium,
  BoldItalicFont=*-MediumItalic,
]{Roboto}

\newcommand*{\footerfont}{\sourcesanspro}
\newcommand*{\bodyfont}{\sourcesanspro}
\newcommand*{\bodyfontlight}{\sourcesansprolight}

%-------------------------------------------------------------------------------
%                Configuration for styles
%-------------------------------------------------------------------------------
\renewcommand{\refname}{}
%-------------------------------------------------------------------------------


%-------------------------------------------------------------------------------
%                Configuration for styles
%-------------------------------------------------------------------------------
% Configure styles for each CV elements
% For fundamental structures



% Boolean value to switch section color highlighting
\newbool{acvSectionColorHighlight}
\setbool{acvSectionColorHighlight}{true}

% Awesome section color
\def\@sectioncolor#1#2#3{%
  \ifbool{acvSectionColorHighlight}{{\color{blue}#1#2#3}}{#1#2#3}%
}


% Colors 
\definecolor{darktext}{HTML}{414141}
\definecolor{darkgray}{HTML}{333333}
\definecolor{gray}{HTML}{5D5D5D}
\definecolor{lightgray}{HTML}{999999}

% Text colors
\colorlet{text}{darkgray}
\colorlet{graytext}{gray}
\colorlet{lighttext}{lightgray}








\newcommand*{\headerfirstnamestyle}[1]{{\fontsize{32pt}{1em}\headerfontlight\color{graytext} #1}}
\newcommand*{\headerlastnamestyle}[1]{{\fontsize{32pt}{1em}\headerfont\bfseries\color{text} #1}}
\newcommand*{\headerpositionstyle}[1]{{\fontsize{7.6pt}{1em}\bodyfont\scshape\color{awesome} #1}}

% \newcommand*{\headeraddressstyle}[1]{{\fontsize{8pt}{1em}\headerfont\itshape\color{lighttext} #1}}

\newcommand*{\headeraddressstyle}[1]{{\fontsize{8pt}{1em}\headerfont\color{lighttext} #1}}
\newcommand*{\headersocialstyle}[1]{{\fontsize{6.8pt}{1em}\headerfont\color{text} #1}}
\newcommand*{\headerquotestyle}[1]{{\fontsize{9pt}{1em}\bodyfont\itshape\color{darktext} #1}}
% footer text style
\newcommand*{\footerstyle}[1]{{\fontsize{8pt}{1em}\footerfont\scshape\color{lighttext} #1}}

\renewcommand*{\sectionstyle}[1]{{\fontsize{12pt}{1em}\bodyfont\bfseries\color{custom} #1}}

\renewcommand*{\subsectionstyle}[1]{{\fontsize{12pt}{1em}\bodyfont\scshape\textcolor{text}{#1}}}

\newcommand*{\paragraphstyle}{\fontsize{9pt}{1em}\bodyfontlight\upshape\color{text}}

% For elements of entry
\newcommand*{\entrytitlestyle}[1]{{\fontsize{10pt}{1em}\bodyfont\bfseries\color{darktext} #1}}
\newcommand*{\entrypositionstyle}[1]{{\fontsize{8pt}{1em}\bodyfont\scshape\color{graytext} #1}}
\newcommand*{\entrydatestyle}[1]{{\fontsize{8pt}{1em}\bodyfontlight\slshape\color{graytext} #1}}
\newcommand*{\entrylocationstyle}[1]{{\fontsize{9pt}{1em}\bodyfontlight\slshape\color{custom} #1}}
\newcommand*{\descriptionstyle}[1]{{\fontsize{9pt}{1em}\bodyfontlight\upshape\color{text} #1}}

% For elements of subentry
\newcommand*{\subentrytitlestyle}[1]{{\fontsize{8pt}{1em}\bodyfont\mdseries\color{graytext} #1}}
\newcommand*{\subentrypositionstyle}[1]{{\fontsize{7pt}{1em}\bodyfont\scshape\color{graytext} #1}}
\newcommand*{\subentrydatestyle}[1]{{\fontsize{7pt}{1em}\bodyfontlight\slshape\color{graytext} #1}}
\newcommand*{\subentrylocationstyle}[1]{{\fontsize{7pt}{1em}\bodyfontlight\slshape\color{custom} #1}}
\newcommand*{\subdescriptionstyle}[1]{{\fontsize{8pt}{1em}\bodyfontlight\upshape\color{text} #1}}

% For elements of honor
\newcommand*{\honortitlestyle}[1]{{\fontsize{9pt}{1em}\bodyfont\color{graytext} #1}}
\newcommand*{\honorpositionstyle}[1]{{\fontsize{9pt}{1em}\bodyfont\bfseries\color{darktext} #1}}
\newcommand*{\honordatestyle}[1]{{\fontsize{9pt}{1em}\bodyfont\color{graytext} #1}}
\newcommand*{\honorlocationstyle}[1]{{\fontsize{9pt}{1em}\bodyfontlight\slshape\color{custom} #1}}

% For elements of skill
\newcommand*{\skilltypestyle}[1]{{\fontsize{10pt}{1em}\bodyfont\bfseries\color{darktext} #1}}
\newcommand*{\skillsetstyle}[1]{{\fontsize{9pt}{1em}\bodyfontlight\color{text} #1}}

\newcommand*{\skilldetailstyle}[1]{{\fontsize{9pt}{0em}\bodyfontlight\color{text} #1}}

% For elements of the cover letter
\newcommand*{\lettersectionstyle}[1]{{\fontsize{14pt}{1em}\bodyfont\bfseries\color{text}\@sectioncolor #1}}
\newcommand*{\recipientaddressstyle}[1]{{\fontsize{9pt}{1em}\bodyfont\scshape\color{graytext} #1}}
\newcommand*{\recipienttitlestyle}[1]{{\fontsize{11pt}{1em}\bodyfont\bfseries\color{darktext} #1}}
\newcommand*{\lettertitlestyle}[1]{{\fontsize{10pt}{1em}\bodyfontlight\bfseries\color{darktext} \underline{#1}}}
\newcommand*{\letterdatestyle}[1]{{\fontsize{9pt}{1em}\bodyfontlight\slshape\color{graytext} #1}}
\newcommand*{\lettertextstyle}{\fontsize{10pt}{1.4em}\bodyfontlight\upshape\color{graytext}}
\newcommand*{\letternamestyle}[1]{{\fontsize{10pt}{1em}\bodyfont\bfseries\color{darktext} #1}}
\newcommand*{\letterenclosurestyle}[1]{{\fontsize{10pt}{1em}\bodyfontlight\slshape\color{lighttext} #1}}




% Define a footer for CV
% Usage: \makecvfooter{<left>}{<center>}{<right>}
\newcommand*{\makecvfooter}[3]{%
  \fancyfoot{}
  \fancyfoot[L]{\footerstyle{#1}}
  \fancyfoot[C]{\footerstyle{#2}}
  \fancyfoot[R]{\footerstyle{#3}}
}

\fancypagestyle{plain}{
  \fancyfoot{} % clear all header and footer fields
  \fancyfoot[L]{December, 2023}
  \fancyfoot[C]{David Sinden~~~·~~~Research Statement}
  \fancyfoot[R]{\thepage\ of \pageref{LastPage}}
}

% document language
\usepackage[english]{babel}  % FIXME: using spanish breaks moderncv

\usepackage[normalem]{ulem}
\usepackage{textcomp}
\usepackage{csquotes}

\newcommand{\myname}[1]{\uline{#1}}

% personal data
\name{David}{Sinden}
\title{Applied Mathematician --- Curriculum Vitae}     % optional, remove / comment the line if not wanted

\email{david.sinden@gmail.com}                         % optional, remove / comment the line if not wanted
\homepage{djps.github.io}                              % optional, remove / comment the line if not wanted

% Social icons
\social[linkedin]{sindendavid}                         % optional, remove / comment the line if not wanted
\social[twitter]{david\_sinden}                        % optional, remove / comment the line if not wanted
\social[github]{djps}                                  % optional, remove / comment the line if not wanted
\social[orcid]{0000-0002-8514-8279}                    % optional, remove / comment the line if not wanted
\social[researchgate]{David-Sinden}                    % optional, remove / comment the line if not wanted
\social[googlescholar]{lauJqakAAAAJ}                   % optional, remove / comment the line if not wanted




%-------------------------------------------------------------------------------
%                letter head definition
%-------------------------------------------------------------------------------
% lengths
%\renewcommand*{\recomputeletterheadlengths}{}

\renewcommand*{\addressfont}{\small\mdseries}

\extrainfo{\href{https://www.github.com/djps}{\githubsocialsymbol}~djps~{\upshape{}|}~\href{https://www.twitter.com/david_sinden}{\twittersocialsymbol}~%
           david\_{}sinden}

% commands
\makeatletter
\newcommand*{\makemyletterhead}{%
  % recompute lengths (in case we are switching from letter to resume, or vice versa)
  \recomputeletterlengths%
  % sender contact info
  \hfill%
  \begin{minipage}{.5\textwidth}%
    % optional detailed information
    \if@details%
      \raggedleft%
      \addressfont\textcolor{color2}{%
        {\bfseries\upshape\@firstname~\@lastname}\@firstdetailselementfalse%
        % optional detailed information
        \ifthenelse{\isundefined{\@addressstreet}}{}{\makenewline\addresssymbol\@addressstreet%
          \ifthenelse{\equal{\@addresscity}{}}{}{\makenewline\@addresscity}% if \addresstreet is defined, \addresscity and addresscountry will always be defined but could be empty
          \ifthenelse{\equal{\@addresscountry}{}}{}{\makenewline\@addresscountry}}%
        \collectionloop{phones}{% the key holds the phone type (=symbol command prefix), the item holds the number
          \makenewline\csname\collectionloopkey \phonesymbol\endcsname\collectionloopitem}%
        \ifthenelse{\isundefined{\@email}}{}{\makenewline\emailsymbol\emaillink{\@email}}%
        \ifthenelse{\isundefined{\@homepage}}{}{\makenewline\homepagesymbol\httpslink{\@homepage}}%
        %\ifthenelse{\isundefined{\@social[twitter]}}{}{\makenewline\twittersocialsymbol\httpslink{\@twitter}}%
        \ifthenelse{\isundefined{\@extrainfo}}{}{\makenewline\@extrainfo}}\fi%
    \end{minipage} \\[2\baselineskip]
  % recipient block
  \begin{minipage}[t]{.5\textwidth}
    \raggedright%
    \addressfont%
    {\bfseries\upshape\@recipientname}\\%
    \@recipientaddress%
  \end{minipage}
  % date
  \hfill% US style
%  \\[1em]% UK style
  \@date\\[2\baselineskip]% US informal style: "January 1, 1900"; UK formal style: "01/01/1900"
  % optional subject
  \ifthenelse{\isundefined{\@subject}}{}{{\bfseries\@subject\\[2\baselineskip]}}
  % opening
  {\centering%
  \textcolor{color1}{\@opening}\\[2\baselineskip]}%
  % ensure no extra spacing after \makelettertitle due to a possible blank line
%  \ignorespacesafterend% not working
  %\hspace{0pt}\par\vspace{-\baselineskip}\vspace{-\parskip}
  }
\makeatother




%----------------------------------------------------------------------------------
%            content
% chrome-extension://efaidnbmnnnibpcajpcglclefindmkaj/https://www.ams.org/journals/notices/202108/rnoti-p1321.pdf?trk=2328&cat=career
%----------------------------------------------------------------------------------


\usepackage{ifxetex}
\ifxetex

  % use support for setting fonts
  \usepackage{fontspec}
  \usepackage[no-sscript]{xltxtra}
  \usepackage{xunicode}

  % allow TeX short cuts such as -- or ``
  \defaultfontfeatures{Mapping=tex-text}

  % set Frutiger as sans serif font (overwrites arev)
  \setsansfont[BoldFont={Frutiger LT Com 65 Bold}, 
               ItalicFont={Frutiger LT Com 56 Italic}, 
               BoldItalicFont={Frutiger LT Com 66 Bold Italic}]{Frutiger LT Com 55 Roman}

  % define really small font for footline
  \font\rsf="Frutiger LT Com 55 Roman" at 4pt

  \def\copyrightsymbol{\copyright}

\else

  \RequirePackage[utf8]{inputenc}
  \usepackage[utf8]{inputenc}
  \def\rsf{\fontsize{4pt}{4pt}\selectfont}
  \def\XeLaTeX{XeLaTeX}
  \def\copyrightsymbol{\textcircled{\hspace{-0.1pt}c}}%Arev does not provide © symbol, but this looks close enough
  \DeclareUnicodeCharacter{00B5}{{\fontfamily{phv}\selectfont \textmu}}%Arev does not have µ symbol, but Helvetica's looks close enough

\fi

\newfontfamily\frutigerfont{Frutiger LT Com 65 Bold}  % or your system’s Frutiger name
\tikzset{
  frutiger node/.style={font=\frutigerfont}
}



\usepackage[sectionbib, square, sort, comma, numbers]{natbib}
\usepackage{chapterbib}

% footer details
\usepackage{lastpage}

\usepackage{bibtopic} %bib by section

\usepackage{eso-pic} % AddToShipoutPictureBG

\begin{document} 

\hypersetup{%
  pdftitle={David Sinden - Project Proposal},
  pdfauthor={David Sinden},
  pdfsubject={Klaus Tschira Stiftung Proposal},
  pdfkeywords={Ultrasound}
}

\AddToShipoutPictureBG*{%
  \begin{tikzpicture}[remember picture, overlay]
    \node[
      anchor=north west,
      inner sep=0pt,
      xshift=0cm,
      yshift=0cm
    ] at (current page.north west) {%
      \includegraphics[width=\paperwidth,  
                       trim=0cm 15cm 0cm 0cm,
                       clip]{fhg_smooth.pdf}%
    };
  \end{tikzpicture}%
}

\makecvfooter{October, 2025}{David Sinden~~~$\cdot$~~~Project Proposal}{\thepage\ ~of~ \pageref{LastPage}}

\noindent\hfill
\scalebox{2}{%
\begin{tikzpicture}[remember picture, overlay]
% \draw ([yshift=-\baselineskip]current page.north) node[inner sep=0, outer sep=0] (icon) { \includegraphics[width=1.0\baselineskip]{fhgicon}};

  \node[
    inner sep=0, outer sep=0, anchor=north west
  ] (icon)
    at ([yshift=-4.5\baselineskip, xshift=-10\baselineskip]current page.north east)
    {\includegraphics[width=1.0\baselineskip]{fhgicon}};

%\node[frutiger node, anchor=south west, xshift=-0.1\baselineskip, yshift=-0.3\baselineskip] (fhg) at (icon.south east) {\Large Fraunhofer};

  \node[
    frutiger node,
    %draw,                 
    %thick,                
    %rounded corners,      
    inner sep=2pt,        
    anchor=south west,
    xshift=-0.1\baselineskip,
    yshift=-0.18\baselineskip
  ] (fhg) at (icon.south east) {\large Fraunhofer};

\node[frutiger node, anchor=north east, yshift=0.4\baselineskip, xshift=0.15\baselineskip] (mevis) at (fhg.south east) {\tiny {MEVIS}}; 
% \node[draw,rounded corners,inner sep=4pt,fit=(icon)(fhg)(mevis)] {};
\end{tikzpicture}
}
%}
%}
%\end{flushright}

\vspace{3cm}

{\centering\Huge\color{FhGreen}Here comes the title of the dissertation}

\vspace{3cm}

\begin{tikzpicture}
  % A node of fixed width, filled with yellow!30
  \node[
    fill=FhGreen!10,        % background color
    % draw=orange,           % optional border
    % thick,                 % border width
    % rounded corners,       % nice corners
    inner sep=1em,         % padding
    text width=\textwidth, % wrap lines at 80% of line width
    align=justify          % or left, center, right
  ] {Controlling the oscillations of microbubbles to safely and reversibly open the blood-brain-barrier is a problem. There little state of the art of the modelling. Koopman operators have been applied to single free bubbles. Here we propose to exploit non-spherical oscillations. };
\end{tikzpicture}

\vfill

\large\color{black}
\begin{tabularx}{\textwidth}{Xr}
& \textbf{\Large David Sinden, PhD}\vspace{0.8cm}\\
%Institute : 
                                                & Microbubbles \\
Fraunhofer Institute for Digital Medicine MEVIS & Drug Delivery \\
                                                & Control Theory %\\[2cm]
%& Submitted 2025
\end{tabularx}

%\vspace{1.5cm}
%\noindent\makebox[\linewidth]{\color{blueXIIlight}\rule{30cm}{8pt}}

\newpage
\thispagestyle{fancy}
\section{Project Idea}

\bibliographystyle{apsrev4-1long} 
\begin{btUnit}
\begin{btSect}[apsrev4-1long]{bibliography} 

I am an interdisciplinary applied mathematician, having worked worked with clinicians, engineers and measurement scientist. I received my Ph.D. in 2009 under the guidance of Prof. Gert van der Heijden at the Centre for Nonlinear Dynamics at University College London, motivated by the static post-buckling configuration of an charged elastic rod in a magnetic field. After my Ph.D. I took a post doctoral position at the department of Mechanical Engineering, also at UCL, as part of the ultrasound research group. The research moved from Hamiltonian dynamical systems to collective behaviour of driven nonlinear oscillators, investigating the oscillations of bubbles which can form in tissue when exposed to high-intensity therapeutic ultrasound.

This lead to an another post-doctoral position at the Institute of Cancer Research, where I designed a treatment planning platform for transcostal high-intensity focused ultrasound. This involved optimization of linear partial differential equations for acoustics on large domains, with constrained on where the acoustic field should be focused, while minimizing damage to surrounding structures. I then obtained a permanent position at the National Physical Laboratory, which is the government laboratory in the United Kingdom. My main responsibility was to lead the standardisation efforts for therapeutic ultrasound. A significant proportion of my work related to validating measurement-based simulations, i.e. simulations on domains computed from imaging data, with initial conditions taken measurement data. I am an active member the IEC technical committee on ultrasound, working to bring the knowledge to a standard. 

At Fraunhofer MEVIS I have worked on a number of projects, expanding expertise into other areas of diagnostic ultrasound, image-guided therapy as well as systems biology. I have recently been awarded a Fraunhofer DISCOVER grant to investigate applications of computational topology in medical imaging.

\section{Computational Topology in Medical Imaging}

Classical medical image analysis is based upon geometric measures, such as lengths, angles, shapes etc. This is somewhat limited as such measures do not fully capture the complexity of information across lengths scales (as these are measured at single spatial scale), or global properties of data. Radiomics seeks to extend classical image analysis by looking at local variations in pixel-to-pixel intensity, giving a measure of texture. Such local measures seek to quantify intuitive concepts such as smoothness or roughness and have been applied to image segmentation. However, the are highly dependent on the imaging modality and are not reproducible. 

Computational topology, and specifically a tool called persistent homology, has recently emerged as powerful tool which can provide insight which is distinct from pixel-based approaches by describing measures of connectedness between objects in an image~. It is referred to characterising the shape of data.  It has also been significant as embedding in highly connected networks. 

I plan to apply techniques from computational topology to two fields: \textcolor{custom}{ultrasound thermometry} and \textcolor{custom}{characterisation of aspects of liver function}, via structure. The first application is to use specular information to correlate persistence diagrams with relative changes in temperature. 

The second application has a number of aspects. The main aim is to quantitatively model changes in perfusion after resection, via persistence diagrams. Another topic will be to develop tools which measure cell alignment, which characterise regeneration and also, from microscopy data, characterise the fenestrations along the liver endothelial cells which influence metabolic rates.

\newpage
Continued~\cite{ravishankar2023sonosamtrack} here.

\newpage
\section{Bibliography}

\btPrintCited 

\end{btSect} 
\end{btUnit} 


\newpage
\section{Contribution}

How does the Boost Fund contribute to your project goals?  

Describe the \enquote{Boost} to your career: Let us know how the funding, career development program and fellowship community make a difference in your career (e.g. for your independence, your development as a researcher, for networking, collaborations, data collection, etc.).  

What potential would remain untapped without this funding? You are welcome to be specific, personal, and authentic. 


\newpage
\section{Time Scales}

Outline the time frame for your project, including preliminary work, intermediate and final project goals. 

What are the key milestones of the project during the funding period?

    \begin{tikzpicture}
      \GanttHeader{\textwidth}{2ex}{4cm}{7}
      \Task{1}{First task}{0.5}{2.5}
      \Task{2}{Second task}{1.5}{2}
      \Task{3}{Another one}{3}{1}
      \Task{4}{Another one}{4}{1}
      \Task{5}{Yet another one}{5}{2}
    \end{tikzpicture}

\end{document}


